% Sign Families 373-417
\section{Sign Families: Complex Forms (M373--M417)}

\subsection{Anthropomorphic Signs (M373--M390)}

Human and humanoid figures represent institutional or role categories:

\subsubsection{Standing Figures}
\begin{itemize}
    \item M373--M378: Standing human forms with variations
    \item M379--M382: Figures with headdresses or crowns
    \item M383--M385: Figures with tools or implements
\end{itemize}

\textbf{Operational Function}: Role specifiers or institutional authority markers.

\subsection{Animal Signs (M391--M410)}

Zoomorphic signs beyond fish:

\subsubsection{Major Animal Categories}
\begin{itemize}
    \item M391--M395: Bovine forms
    \item M396--M400: Other quadrupeds
    \item M401--M405: Birds
    \item M406--M410: Reptiles and other fauna
\end{itemize}

\textbf{Operational Function}: Domain specifiers, institutional emblems, or commodity classifiers.

\subsection{Terminal Signs (M411--M417)}

Signs appearing primarily at sequence endings:

\subsubsection{Common Terminals}
\begin{itemize}
    \item M411: Enclosure sign
    \item M412--M414: Roof-like forms
    \item M415--M417: Abstract closures
\end{itemize}

\textbf{Operational Function}: Closure operators, validation markers, or seal signatures.

\subsection{Complex Co-occurrence Patterns}

Advanced combinatorial analysis reveals:

\begin{itemize}
    \item Anthropomorphic signs precede animal signs in ordered sequences
    \item Animal signs can combine with counting strokes (quantity specification)
    \item Terminal signs rarely co-occur (mutual exclusion)
    \item Certain complex forms never appear together (operational incompatibility)
\end{itemize}

\subsection{Variant Relationships}

Many signs in this range show systematic variations:
\begin{itemize}
    \item Orientation changes (left-facing vs. right-facing)
    \item Elaboration levels (simple vs. ornamented)
    \item Compositional structure (monolithic vs. compound)
\end{itemize}

These variants often preserve operational function while marking contextual or stylistic differences.
